\documentclass[]{asc}
\usepackage[final]{graphicx}
\usepackage{amsmath}
\usepackage{tikz}
\usepackage[english]{babel}
\usepackage{verbatim}
\usetikzlibrary{arrows,shapes,positioning}
\usepackage[colorlinks,%
            linkcolor=black,%
            citecolor=black,%
            filecolor=black,%
            urlcolor=black]{hyperref}
\usepackage{subfigure}
\usepackage[hang,flushmargin]{footmisc}
\captionsetup[figure]{name={Figure},labelsep=period}
\captionsetup[table]{labelsep=period}

\makeatletter
\renewcommand\@biblabel[1]{#1.}
\makeatother


\begin{document}


\title{Paper Guidelines for Proceedings of the American Society for Composites 32nd Technical Conference}
\author{\bf Anthony Deraco\\ \bf Joseph Eckenrode \\ Paper Number: XXX}

%\affiliation{American Society for Composites}

\maketitle

\begin{abstract}

Below is a set of instructions for preparing your paper in a format suitable for the intended media for these proceedings. It is most important that you follow these guidelines as closely as possible and adhere to the margins (discussed below) for overall consistency and format relationship to all other papers in the intended proceedings. Print or type text as close to the margins as possible without going beyond these margins. Abstract should not exceed 300 words.



\end{abstract}

\let\thefootnote\relax\footnote{
\fontfamily{phv}\fontsize{10}{12}\selectfont
Anthony Deraco,and Joseph Eckenrode \\ American Society for Composites}

\section*{PAPER SUBMISSION}

\textcolor{blue}{\it All the information of this section can be deleted}

PAPER DEADLINE: June 16, 2017

PAPER LENGTH: 6-16 PAGES (Maximum, not counting cover page)\\ 


SUBMISSION PROCEDURE: Information on the electronic submission of
manuscripts is provided on the conference web site - 
\textbf{https://cdmhub.org/groups/asc2017}\\

FILE NAME: \texttt{PaperNo\_LastName(1st author).pdf}

EXAMPLE: \texttt{1500\_Deraco.pdf}\\

Please be sure to read these LaTeX Guidelines prior to preparing your paper. This will ensure your paper is consistent with the format of the articles in the CD-ROM. \\

\textbf {NOTE: If you are using LaTeX to prepare your manuscript, please submit your paper as PDF file. Sample guidelines are shown with the correct margins. Follow the style from these guidelines for your page format.}\\

NOTES
\begin{itemize}
\item Use this LaTeX package as a template to prepare your paper  
\item First page will be the Cover Page
\item Do not include page numbers
\item File size should not be greater than 10 Mb. Resize/compress images as needed.
\end{itemize}


\section*{PRINTING/FORMATTING YOUR MANUSCRIPT}

The margins to be used for your paper are as follows: For Top 0.75'' (19.05 mm); Bottom 0.93'' (23.6 mm); Left 1.32'' (33.5 mm); Right 1.56'' (39.6 mm). In Page Layout, select Size and use LETTER (8.5 x 11 inches). Make sure that the width of the text is a maximum of 5.625 inches (143 mm) and the depth is a maximum of 9.25 inches (235 mm). Adjust indents if using A4 size.
Use Times New Roman (the typeface used in these instructions), 12 point, at twelve characters per inch or in proportional spacing (as here).

\section*{SUBHEADS}\textcolor{blue}{\it One line space from heading to text} 

First-order subheads should be in bold caps with 2 line spaces above and 1 line space below them. Second-order subheads should be in bold with main words capitalized with 1 line space above and below them. Third-order subheads should be in regular type in all caps with 1 line space above and below them. Subheads should not appear alone at the bottom of a page.


\subsection*{Subsection}\textcolor{blue}{\it Second-order heading}

Do not leave a blank line between paragraphs. The first line of each paragraph should be indented 18 points (1.5 pica-6.3 mm-1/4'') space.
(two line spaces between text and first-order subhead)
 


\section*{FIRST PAGE}
DEStech Publications, Inc. will insert the title and author name(s) in the space left blank on the top of the first page. The title and author name(s) should be supplied on a separate sheet (cover page) and submitted with the manuscript. {\bf The author's current affiliation, including full address, should appear as a footnote on the bottom of the first page in a 10-point Helvetica/Univers/Arial font. If authors are at different affiliations, please include names with affiliations or superscript numbers {\it (see cover page)}}. A 1-point rule (1 inch long, 25.4 mm) should be no more than $1/8$ inch (3 mm) above the affiliation. A SUMMARY should begin the paper three and a one  quarter inches (83 mm) down from the top (to leave space for the title and author name). \underline{No illustrations, figures, or tables should appear on the first page}.


\section*{ARTWORK}

Artwork (photographs, tables, or figures) should be planned for appearance at the bottom or top of a page, as close to the first mention of the artwork as possible; i.e., no text should appear both above and below artwork. There should be between two and six blank lines between artwork (including captions) and main text and one blank line between artwork and caption (see Figure~\ref{fi.FigureBox}). Artwork should be centered between the left and right margins. In cases where artwork takes an entire page, it should be centered within the margins. All print within the artwork should be in Helvetica/Univers/Arial font in the range of 8 point to 10 point (see Figure~\ref{fi.ExampleFigure}).

\subsection*{Tables and Figures}

All tables and figures with text only should be boxed in; i.e., a box should be drawn around the table or figure either by hand with a ruler or with a draw facility on the computer (see Figure~\ref{fi.FigureBox}  sample and Table \ref{ta.ExampleTable} sample). Labels for tables should appear centered at the top of the table in uppercase letters in a 10-point font. \textbf{A Roman numeral should follow the word ``TABLE''. Tables should be numbered consecutively throughout the text, i.e. I, II, III, IV, etc. In the case of figures (charts, graphs, photographs), the word ``Figure'' should appear preceding an Arabic numeral (1, 2, 3, etc.) in the caption, which should be centered below the figure}. Figures should also be numbered consecutively throughout the text. Captions must always be printed or typed, not handwritten.

\textcolor{blue}{\it Two to six blank lines , text figures ``boxed in''}
\newline
\newline
\begin{figure}[b]
\noindent\fbox{%
	\parbox{\textwidth}{%
		\begin{itemize}
			\item \textbf{FIRST-ORDER HEADINGS (e.g. ABSTRACT):} bold uppercase; 2 line spaces above heading and 1 line space below heading.
			\item \textbf{Second-order Headings (e.g. Tables and Figures):} bold; initial cap of main words; 1 line space above and below heading.
			\item  THIRD-ORDER HEADINGS (e.g. PHOTOGRAPHS): bold; regular font; all caps; 1 line space above and below heading.
		\end{itemize}
	}%
}
\caption{Subhead specifications, (example of a text figure, text figures ``boxed in'')}
\label{fi.FigureBox}
\end{figure}

\begin{figure}[t]
 \centering
 \includegraphics[width=250pt]{ExampleOfFigure.png}
 \caption{Example of the figure.}
 \label{fi.ExampleFigure}
\end{figure}

\textcolor{blue}{\it Two to six blank lines from figure to text or text to figure}
\newline
\newline 
\subsubsection*{PHOTOGRAPHS} \textcolor{blue}{\it Third-order heading}

Original color photographs are preferred (black and white is acceptable if color is not available) and should be submitted with your manuscript if you do not insert halftones into your article. \underline{Inserting a 300 dpi halftone is the preferred method}. If photographs will be submitted separate from the main text, they should be clearly marked as to where within the text they belong.

\subsubsection*{ILLUSTRATIONS} 

\textcolor{blue}{\it Third-order heading}

Illustrations should be treated as ``Figures'' and labeled as such. Illustrations should be professionally drawn using black India ink (do not use blue ink) or created with a draw facility on the computer. Labels must always be printed or typed, not handwritten.

\textcolor{blue}{\it Two to six blank lines from table to text or text to table.}
\newline
\newline 
\begin{table}[t]
	\centering
	\caption{experimental and calculated data}
	\begin{tabular}{|cccc|}
		\hline
		$V_{f}$ & $\alpha I$ & $\Delta T$ & $\Delta \theta/T$\\ 
	    $(\%)$ & $\xi 10^{-6}/F$ & (F)& x $10^{-3}/F$\\
		\hline
		55 &  1.1 & -150 &  -1.90\\
		35T1.5 & -235 & -2.55 &  28.74\\
		35 & 1.50&-235&-2.40\\
		60 &  0.8&-235&-1.60\\
		60 & 0.8&-235&-1.65\\
	    60 & 0.8 &-235&-1.50\\
		\hline
	\end{tabular}
	\label{ta.ExampleTable}
\end{table}

\section*{MATHEMATICAL EQUATIONS} 
Mathematical equations should be centered between left and right margins and should be separated by one blank line. All equations should be numbered consecutively throughout the text, using Arabic numerals in parentheses along the right margin.

\begin{equation}
i = 1 
\label{eq:OneExample}
\end{equation}

\begin{equation}
a + b = c 
\label{eq:TwoExample}
\end{equation}

\section*{REQUIREMENTS FOR REFERENCES} 
The first time a particular source (article, book, etc.) is cited in the text, it should be assigned a number and placed in the list of references at the end of the paper, similar to the list that follows. Each time this source is cited in the text, it should be referred to by its original number. Each number should be placed in brackets at the end of the referenced phrase before any final punctuation mark, e.g., as [2]. A sample of listings is provided for a journal paper \cite{Ikegami}, a book \cite{Mitsiti}, a chapter in an edited book \cite{Inman}, an unpublished report \cite{Margarit}, and a presentation at a conference \cite{Hoffer}. (If square brackets are not available to you, parentheses are acceptable.) \textbf{References should be set in 10-point type.}


\begin{thebibliography}{9}
\bibitem{Ikegami}
Ikegami, R., D. G. Wilson, J. R. Anderson, and G. J. Julien. 1990. ``Active Vibration Control Using NiTiNOL and Piezoelectric Ceramics," {\it J. Intell. Matls. Sys. \& Struct.}, 20(2):189-206.
\bibitem{Mitsiti}
Mitsiti, M. 1996. {\it Wavelet Toolbox, For Use with MALAB}. The Math Works, Inc., pp. 111-117.
\bibitem{Inman}
Inman, D.J. 1998. ``Smart Structures Solutions to Vibration Problems,"  in {\it International Conference on Noise and Vibration Engineering}, C. W. Jefford, K. L. Reinhart, and L. S. Shield, eds. Amsterdam: Elsevier, pp. 79-83.
\bibitem{Margarit}
Margarit, K. L. and F. Y. Sanford. March 1993. ``Basic Technology of Intelligent Systems," Fourth Progress Report, Department of Smart Materials, Virginia Polytechnic Institute and State University, Blacksburg.
\bibitem{Hoffer}
Hoffer, R. and D. Dean. 1996. ``Geomatics at Colorado State University," presented at the 6th Forest Service Remote Sensing Applications Conference, April 29-May 3, 1996.	
\end{thebibliography}

\noindent\fbox{%
	\parbox{\textwidth}{%		
\section*{IMPORTANT NOTES} 
\begin{enumerate}
\item Your paper should be formatted single-spaced, and you should adhere strictly to the margins outlined in the section on ``Printing/Typing Your Manuscript." Stay within all margin lines, but fill the page as closely to the margins as possible.
\item DO NOT use San Serif font or Italic faces for the main body of the text.
\item Please supply the electronic file as a PDF.
\item If a graph or table must be turned sideways for viewing, please be certain that the top of the table is on the left of the page. There also should not be any text (besides figure or table caption) on a page that has a graph or table turned sideways.
\item Questions regarding format should be addressed to Tony Deraco or Steve Spangler, DEStech Publications, Inc.; tel.: 717/290-1660; fax: 717/509-6100; e-mail: \underline{aderaco@destechpub.com} or \underline{sspangler@destechpub.com}.
\item For content questions, contact the committee chairperson or the person to whom you should forward Inquiries to. 
\end{enumerate}
	}%
}
\end{document}
